\documentclass[11pt, a4paper]{article}

%%% SST LAB PROTOCOLL PREAMBLE
%%% 2019
%%%%%%%%%%%%%%%%%%%%%%%%%%%%%%%


%%% PACKAGES
%%%%%%%%%%%%%%%%%%%%%%%%%%%

\usepackage[ngerman]{babel}

\usepackage[utf8]{inputenc}
\usepackage{amsmath}
\usepackage{pgfplots}
\usepackage{tikz}
\usepackage[many]{tcolorbox}
\usepackage{graphicx}
\graphicspath{ {./graphics/} }
\usepackage{pdfpages}
\usepackage{dashrule}
\usepackage{float}
\usepackage{siunitx}
\usepackage{trfsigns}
\usepackage{booktabs}
\usepackage[european]{circuitikz}
\usepackage{tcolorbox}

%%% DOCUMENT GEOMETRY
%%%%%%%%%%%%%%%%%%%%%%%%%%%

\usepackage{geometry}
\geometry{
 a4paper,
 total={0.6180339887498948\paperwidth,0.6180339887498948\paperheight},
 top = 0.1458980337503154\paperheight,
 bottom = 0.1458980337503154\paperheight
 }
\setlength{\jot}{0.013155617496424828\paperheight}
\linespread{1.1458980337503154}

\setlength{\parskip}{0.013155617496424828\paperheight} % paragraph spacing


%%% COLORS
%%%%%%%%%%%%%%%%%%%%%%%%%%%

\definecolor{red1}{HTML}{f38181}
\definecolor{yellow1}{HTML}{fce38a}
\definecolor{green1}{HTML}{95e1d3}
\definecolor{blue1}{HTML}{66bfbf}
\definecolor{hsblue}{HTML}{00b1db}
\definecolor{hsgrey}{HTML}{afafaf}

%%% CONSTANTS
%%%%%%%%%%%%%%%%%%%%%%%%%%%
\newlength{\smallvert}
\setlength{\smallvert}{0.0131556\paperheight}


%%% COMMANDS
%%%%%%%%%%%%%%%%%%%%%%%%%%%

% differential d
\newcommand*\dif{\mathop{}\!\mathrm{d}}

% horizontal line
\newcommand{\holine}[1]{
  	\begin{center}
	  	\noindent{\color{hsgrey}\hdashrule[0ex]{#1}{1pt}{3mm}}\\%[0.0131556\paperheight]
  	\end{center}
}

% mini section
\newcommand{\minisec}[1]{ \noindent\underline{\textit {#1} } \\}

% quick function plot
\newcommand{\plotfun}[3]{
  \vspace{0.021286\paperheight}
  \begin{center}
    \begin{tikzpicture}
      \begin{axis}[
        axis x line=center,
        axis y line=center,
        ]
        \addplot[draw=red1][domain=#2:#3]{#1};
      \end{axis}
    \end{tikzpicture}
  \end{center}
}

% box for notes
\newcommand{\notebox}[1]{

\tcbset{colback=white,colframe=green1!100!black,title=Note!,width=0.618\paperwidth,arc=0pt}

 \begin{center}
  \begin{tcolorbox}[]
   #1 
  \end{tcolorbox}
 
 \end{center} 
 
}

% box for equation
\newcommand{\eqbox}[2]{
	
	\tcbset{colback=white,colframe=green1!100!black,title=,width=#2,arc=0pt}
	
	\begin{center}
		\begin{tcolorbox}[ams align*]
				#1
		\end{tcolorbox}
		
	\end{center} 
	
}
% END OF PREAMBLE

\begin{document}

\includepdf{./titlepage/titlepage.pdf}

\section{Teil A: 2-Phasen-Scheduler}

\subsection{Prototyp 2-Phasen-Scheduler sched2p}

\lstinputlisting[title={sched2p.py}]{code/2phase.py}%[language=Python]

\begin{lstlisting}[title={Ausgabe sched2p bei Simulationszeit von 10 Minuten}]
time,state
0.0,TrafficLightStates.RED
180.12345,TrafficLightStates.GREEN
240.69134,TrafficLightStates.RED
420.81479,TrafficLightStates.GREEN
481.38268,TrafficLightStates.RED
\end{lstlisting}

\subsection{Zeitgetriebener Simulator (sched2p\_timedriven)}

\lstinputlisting[title={sched2p\_timedriven.py}]{code/2phase_timedriven.py}%[language=Python]

\begin{lstlisting}[title={Teil der Ausgabe sched2p\_timedriven bei Simulationszeit von 10 Minuten. Schrittweite = $1\,\si{\second}$}]
time,state
...
422.0,TrafficLightStates.RED
423.0,TrafficLightStates.RED
424.0,TrafficLightStates.RED
425.0,TrafficLightStates.GREEN
426.0,TrafficLightStates.GREEN
427.0,TrafficLightStates.GREEN
...
483.0,TrafficLightStates.GREEN
484.0,TrafficLightStates.GREEN
485.0,TrafficLightStates.GREEN
486.0,TrafficLightStates.RED
487.0,TrafficLightStates.RED
488.0,TrafficLightStates.RED
...
599.0,TrafficLightStates.RED
600.0,TrafficLightStates.RED
601.0,TrafficLightStates.RED
terminated
\end{lstlisting}

Die Ausgabe bei $601$ Sekunden liegt daran, dass im Simulationsschritt die Zeit nach dem Ausführen des Ereignisses erhöht wird. Die eigentliche Zeit bei Eintreten des Ereignisses liegt also immer $1\,\si{s}$ früher als es die Ausgabe zeigt.


\subsection{Vergleich von sched2p und sched2p\_timedriven}

Man erkennt gut aus der Ausgabe von \inlinecode{sched2p_timedriven}, dass die Zeitpunkte der Ereigniswechsel nicht genau getroffen werden können und die Genauigkeit immer von der Zeitauflösung / Schrittweite abhängt. Höhere Genauigkeit bringt allerdings auch eine erhöhte Anzahl an Iterationen mit sich, wodurch die Simulationszeit (Echtzeit) verlängert wird.

Eine Ausgabe der Anzahl der Iterationen zeigt $600$, wie erwartet ($\frac{600\,\si{\second}}{1\,\si{\second}/\text{step}}$). Im Gegensatz zum Ereignisbasierten Scheduler mit $5$ Iterationen liegt der Unterschied also bei einem Faktor von $120$.

Außerdem zeigt der zeitbasierte Ansatz (in diesem Fall), dass die meisten Iterationen keine Zustandsänderungen herbeiführen und damit nutzlos sind.

\paragraph{Simulationszeit von $1$ Woche:}
Eine Woche entspricht $604800$ Sekunden. Mittels \inlinecode{time()}-Funktion aus der \inlinecode{time} library in Python konnten die Ausführungszeiten beider Programme geprüft werden.

\[t_{\text{sched2p}} = 0.00652\,\si{\second}\]
\[t_{\text{timedriven}} = 0.182 \,\si{\second}\]
\[\frac{t_{\text{timedriven}}}{t_{\text{sched2p}}} \approx 30\]

\section{Teil B: 3-Phasen Scheduler}
\lstinputlisting[title={sched3p.py}]{code/3phase.py}

\begin{table}[H]
  \centering
\begin{tabular}{|c|c|c|l|}
\hline
\multicolumn{1}{|l|}{Simulationslauf} & \multicolumn{1}{l|}{Rot-Rot-Übergänge} & \multicolumn{1}{l|}{Rot-Grün-Übergänge} & Zustandsverlauf \\ \hline
1                                     & 1                                      & 1                                       & R-R-G-R         \\ \hline
2                                     & 0                                      & 2                                       & R-G-R-G-R       \\ \hline
3                                     & 1                                      & 1                                       & R-G-R-R         \\ \hline
4                                     & 1                                      & 1                                       & R-R-G-R         \\ \hline
5                                     & 2                                      & 0                                       & R-R-R-G         \\ \hline
6                                     & 1                                      & 1                                       & R-R-G-R         \\ \hline
7                                     & 1                                      & 1                                       & R-R-G-R         \\ \hline
8                                     & 3                                      & 0                                       & R-R-R-R         \\ \hline
9                                     & 0                                      & 2                                       & R-G-R-G-R       \\ \hline
10                                    & 1                                      & 1                                       & R-G-R-R         \\ \hline
\end{tabular}
\caption{Simulationsergebnis sched3p bei 10 Druchläufen mit Simulationszeit von $10\,\si{\minute}$}
\end{table}

\begin{table}[H]
\centering
\begin{tabular}{cc|c|c|}
\hline
\multicolumn{1}{|l|}{Simulationslauf} & \multicolumn{1}{l|}{Rot-Rot-Übergänge} & \multicolumn{1}{l|}{Rot-Grün-Übergänge} & \multicolumn{1}{l|}{Verhältnis (R-R/R-G)} \\ \hline
\multicolumn{1}{|c|}{1}               & 5851                                   & 5672                                    & 0.9694                                    \\ \hline
\multicolumn{1}{|c|}{2}               & 5732                                   & 5761                                    & 1.0051                                    \\ \hline
\multicolumn{1}{|c|}{3}               & 5695                                   & 5789                                    & 1.0165                                    \\ \hline
\multicolumn{1}{|c|}{4}               & 5681                                   & 5799                                    & 1.0208                                    \\ \hline
\multicolumn{1}{|c|}{5}               & 5749                                   & 5748                                    & 0.9998                                    \\ \hline
\multicolumn{1}{|c|}{6}               & 5729                                   & 5763                                    & 1.0059                                    \\ \hline
\multicolumn{1}{|c|}{7}               & 5743                                   & 5753                                    & 1.0017                                    \\ \hline
\multicolumn{1}{|c|}{8}               & 5745                                   & 5751                                    & 1.0010                                    \\ \hline
\multicolumn{1}{|c|}{9}               & 5733                                   & 5760                                    & 1.0047                                    \\ \hline
\multicolumn{1}{|c|}{10}              & 5732                                   & 5761                                    & 1.0051                                    \\ \hline
\multicolumn{1}{l}{}                  & \multicolumn{1}{l|}{}                  & Mittelwert                              & 1.003                                     \\ \cline{3-4}
\end{tabular}
\caption{Simulationsergebnis sched3p bei 10 Druchläufen mit Simulationszeit von $4$ Wochen}
\end{table}

\end{document}
